\textbf{Predstavil:} Taj Zupan

\begin{enumerate}
  \item Dana je krožnica $K \subset \mathbb{R}^{2}$ s središčem v izhodišču in polmerom $1$. Za vsako točko $M \in K$ naj bo $P$ pravokotna projekcija $M$ na abscisno os in $Q$ pravokotna projekcija točke $P$ na premico skozi točko $M$ in izhodišče. Določi enačbo krivulje $L$, ki jo opiše točka $Q$ (ko $M$ potuje po krožnici $K$) in jo najbolj natančno skicirajte. 
  \item Krivulja $P$ naj bo dana z enačbo $y^{2} = x$. Za vsako točko $M\in P \setminus \left\{ \left( 0,0 \right) \right\}$ naj bo $A_{M}$ tista točka na poltraku iz točke $\left( 0,0 \right)$ skozi točko $M$, za katero je produkt oddaljenosti točke $M$ od izhodišča in oddaljenosti točke $A_{M}$ od izhodišča enak $1$. Čim bolj natančno skicirajte krivuljo $K=\left\{ A_M \colon M \in P \setminus \left\{ \left( 0,0 \right) \right\} \right\}$.
\end{enumerate}